\documentclass[11pt,a4paper]{article}
\usepackage[margin=24mm]{geometry}
\usepackage{amsmath,amssymb,amsthm,booktabs,longtable,array,hyperref}
\setlength{\emergencystretch}{2em}
\newtheorem{theorem}{Theorem}[section]
\newtheorem{lemma}[theorem]{Lemma}
\newcommand{\Aut}{\operatorname{Aut}}
\newcommand{\Irr}{\operatorname{Irr}}
\newcommand{\ON}{\mathrm{O'N}}
\newcommand{\gen}[1]{\langle #1\rangle}
\title{Regular maps for the O'Nan group:\\a corrected source synthesis}
\author{Human Compression project}
\date{}
\begin{document}
\maketitle
\begin{abstract}
We assemble the acquired proofs and finite certificates for Connor and
Leemans, \emph{Algorithmic enumeration of regular maps}. The corrected
triangular table has 136 entries and sum 798644. Its positive types number
129; the chiral, inner-reflexible and outer-reflexible type spectra have
sizes 128, 95 and 97. We state the counting convention, the imported group
and character data, the source corrections, and the separate evidence for
each of the source's mathematical clauses. The construction and verification
of witnesses are distinguished from exhaustion of markings. In particular,
we do not claim the total numbers of inner- or outer-reflexible markings.
This manuscript joins already acquired components; it does not claim a
new classification of a sporadic group or a new general enumeration theorem.
\end{abstract}

\section{Objects, conventions and the corrected result}\label{sec:scope}
Let $G=\ON$ and $A=\Aut(G)=G:2$. For $p,q>2$ put
\[
 E_G(p,q)=\{(R,S)\in G^2:|R|=p,\ |S|=q,\ |RS|=2,
                         \gen{R,S}=G\},\qquad
 n(p,q)=|E_G(p,q)/A|.
\]
The action is simultaneous application of an automorphism to both entries.
Pairs are ordered. Displaying only $p\leq q$ chooses one orientation of
different-order types; when $p=q$, no additional division by duality is
performed. The source sometimes writes braces around the pair. Throughout
this treatment its numerical convention is the ordered-pair convention
of its equation (2.4). A quotient identifying dual markings is a different
problem.

A marking is reflexible if an automorphism simultaneously inverts $R,S$.
It is inner-reflexible or outer-reflexible according as this automorphism
lies in $G$ or $A\setminus G$; otherwise it is chiral. These alternatives
are disjoint. Their map and polyhedron interpretation uses the
correspondence proved in Section~\ref{sec:strings} and the classical
map interpretation retained from the source.

Write
\begin{align*}
 P&=\{3,4,5,6,7,8,10,11,12,14,15,16,19,20,28,31\},\\
 Z&=\{(3,3),(3,4),(3,5),(3,6),(3,7),(3,12),(4,4)\}.
\end{align*}
and $T=\{(p,q)\in P^2:p\leq q\}\setminus Z$.
\begin{theorem}\label{thm:main}
Under the named inputs in Section~\ref{sec:inputs}, the acquired component
proofs and certificates establish the following corrected conclusions.
\begin{enumerate}
\item The exact values of $n(p,q)$ are Tables~\ref{tab:counts1}--\ref{tab:counts2}.
Their sum is $798644$; their positive support is $T$, of size $129$.
\item Chiral markings occur precisely at $T\setminus\{(3,15)\}$, of size $128$.
\item Inner-reflexible markings occur precisely at
\[
 I=\{(p,q)\in T:20,31\notin\{p,q\}\}
       \setminus\{(3,10),(4,5),(4,6)\},\qquad |I|=95.
\]
\item Outer-reflexible markings occur precisely at
\[
 J=\{(p,q)\in T:8,16\notin\{p,q\}\}\setminus\{(6,12)\},
 \qquad |J|=97.
\]
\end{enumerate}
\end{theorem}
The table counts automorphism orbits of markings. The spectrum sizes
$129,128,95,97$ count \emph{types}. No total of all inner-reflexible
or all outer-reflexible markings is inferred from these type counts.
The proof of the theorem is the following synthesis of the explicitly
listed component proofs, not a new replay of every numerical calculation.

\section{Named inputs and their provenance}\label{sec:inputs}
The following are explicit imported mathematical data, rather than facts
deduced merely by matching a group order or a program label.
\begin{enumerate}
\item The identification of the O'Nan simple group, its order
$g=460815505920$, its centrelessness, its full automorphism group $G:2$,
and its nonidentity element-order set $P\cup\{2\}$ are retained Atlas
inputs. The acquired degree-$245520$ permutation representation of
$\ON:2$ is faithful. Its named socle has an orbit of size $122760$.
The literal permutation readers check words and equalities in this model;
they do not independently classify the named group.
\item The complete subgroup-class roster, normalizer orders, maximal
classes and the specifically used direct inclusion multiplicities are
retained from Connor--Leemans' subgroup-lattice atlas. There are thirteen
maximal classes, of orders
\[
\begin{split}
&3753792,3753792,175560,161280,25920,25920,\\[-2pt]
&14880,14880,10752,7920,7920,2520,2520.
\end{split}
\]
The full roster has 581 subgroup classes. Unweighted containment alone
does not supply arbitrary weighted incidences. Each component requiring
an actual inclusion multiplicity proves or acquires that multiplicity;
it does not multiply weights along arbitrary cover paths.
\item The ordinary character tables named \texttt{ON}, \texttt{ON.2}
and the subgroup tables used in the numerical components are the acquired
Atlas/CTblLib data. Their stored class fusions, order maps, power maps and
identifications are retained inputs. Character arithmetic and class joins
are checked exactly. Column norms and order consistency supplement these
inputs; they are not replacements for their named identifications.
\item The complete centralizer, pair and word-transport certificates
used for involution enumeration are imported at their precise acquired
scope. In particular $|C_G(z)|=161280$ for an inner involution $z$,
$|C_A(z)|=322560$, and an outer involution $t$ has
$C_A(t)=\gen t\times J_1$. The acquired commuting domains have 1750
inner involutions other than $z$, and 1463 outer involutions other than
$t$. Small-action faithfulness and full stabilizers are separately
checked before their orbit counts are used.
\end{enumerate}
The classical Frobenius character identity, Hall subgroup inversion and
the map interpretation retain their original attribution. Woldar's
non-Hurwitz theorem supplies the cited exclusion $(3,7)$; our use of that
theorem does not constitute a recovery of Woldar's entire paper. The
Conder--Oliveros attribution for the simple-group polyhedron interface
is retained, together with the independent intersection-property proof
below. Source reports of old running times and earlier authors' results
remain historical attributions.

The companion relative-path binding list identifies the exact retained
source, model, table, proof and review interfaces. An acceptance record
certifies that the indicated component was checked; it is not substituted
for the component's mathematical proof or for its stated borrowed inputs.

\section{Counting pairs without overcounting subgroups}\label{sec:counting}
For a finite group $H$, let $X_H(p,q)$ be the set defined as $E_H(p,q)$
without the generation condition, and write $x_H=|X_H|$, $e_H=|E_H|$.
For a conjugacy-class representative $c_i$ and an irreducible character
$\chi$, put $v_{\chi,m}=\sum_{|c_i|=m}|c_i^H|\chi(c_i)$.
Frobenius' formula gives
\begin{equation}\label{eq:frob}
 x_H(p,q)=\frac1{|H|}\sum_{\chi\in\Irr(H)}
       \frac{v_{\chi,2}v_{\chi,p}v_{\chi,q}}{\chi(1)}.
\end{equation}
Indeed $(R,S)\mapsto((RS)^{-1},R,S)$ is a bijection to the
order-$(2,p,q)$ triples with product one. The sum over their conjugacy
classes is exactly \eqref{eq:frob}. All arithmetic is in the exact
character field, with the final nonnegative integer checked exactly.

\begin{lemma}\label{lem:partition}
For every finite $H$,
\[
 e_H(p,q)=x_H(p,q)-\sum_{K<H}e_K(p,q).
\]
Grouping actual subgroups by $H$-conjugacy replaces each term by
$[H:N_H(K)]e_K(p,q)$. The action of $\Aut(H)$ on $E_H(p,q)$ is free.
\end{lemma}
\begin{proof}
Every pair in $X_H$ generates one unique subgroup, so the sets $E_K$ for
actual $K\leq H$ partition $X_H$. Conjugate subgroups have equal intrinsic
counts and their number is the normalizer index. An automorphism fixing
both generators fixes every word in them, hence the whole group.
\end{proof}
Consequently the correct denominator for $G$ is $|A|=2g$, applied to
generating pairs, not to all character-counted pairs. In particular,
\[
 n(p,q)=\frac1{2g}\left(x_G(p,q)\ -\
           \#\{\hbox{nongenerating pairs in }X_G(p,q)\}\right).
\]
One may use a complete weighted subgroup matrix and solve a unit
triangular system in increasing subgroup order. The lcm filter is valid
because a subgroup containing elements of orders $p,q$ has order divisible
by $\operatorname{lcm}(p,q)$. This is the source's finite dynamic program
and Hall inversion, not a new inversion theorem.

For this application the numerical components instead exploit exact
maximal-overgroup and generated-subgroup decompositions wherever possible.
A raw union bound over the thirteen maximal classes first gives lower
bounds for $e_G$. It becomes equality only after uniqueness or all
overlaps have been proved. The proofs use direct normalizer incidences,
characteristic subgroups, small literal models, and separate generated
subgroup types to account for these overlaps. This distinction matters:
the bare 581-class roster would not by itself justify all 136 numbers.

The numerical component atlas in the binding list includes the prime
columns $11,19,31$, the $5/7$ and $15$ arguments, the block and normal
overgroup arguments, the complete columns $14,16,20,28$, the remaining
small pairs and the final projective order-six calculation. The latter
joins the exact projective contribution outside involution centralizers
($68g$ and $200g$) to the separately proved complementary contributions
($2g$ and $46g$), giving $(4,6)=43$ and $(6,6)=165$ after subtraction and
division by $2g$. The accepted exact ledger joins all 136 entries to these
proofs and their reviewed input contracts.

\begin{table}[htbp]
\centering\small
\setlength{\tabcolsep}{5pt}
\begin{tabular}{r*{8}{r}}\toprule
$p\backslash q$&3&4&5&6&7&8&10&11\\\midrule
3&0&0&0&0&0&10&7&37\\
4&&0&18&43&102&284&285&503\\
5&&&26&98&150&470&365&718\\
6&&&&165&354&1122&953&1776\\
7&&&&&648&1848&1506&2687\\
8&&&&&&5424&4460&8096\\
10&&&&&&&3613&6532\\
11&&&&&&&&11839\\
\bottomrule\end{tabular}
\caption{Exact ordered $A$-orbit counts, first half of the triangular table.}
\label{tab:counts1}
\end{table}
\begin{table}[htbp]
\centering\small
\setlength{\tabcolsep}{5pt}
\begin{tabular}{r*{8}{r}}\toprule
$p\backslash q$&12&14&15&16&19&20&28&31\\\midrule
3&0&10&6&68&57&44&20&25\\
4&120&166&234&1292&846&554&370&359\\
5&211&290&340&1966&1242&800&560&502\\
6&474&597&874&4700&2943&1948&1370&1268\\
7&815&1054&1284&7448&4725&2916&2214&1995\\
8&2370&3056&3960&22040&13926&8880&6276&5752\\
10&1969&2526&3072&17822&11262&7224&4992&4632\\
11&3583&4601&5814&32488&20493&13094&9202&8325\\
12&1072&1391&1710&9764&6126&3984&2796&2577\\
14&&1796&2330&12504&7899&5024&3616&3266\\
15&&&2834&15808&10020&6424&4496&4062\\
16&&&&88784&56052&35644&25316&23048\\
19&&&&&35442&22572&15978&14553\\
20&&&&&&14238&10292&9118\\
28&&&&&&&7246&6658\\
31&&&&&&&&5999\\
\bottomrule\end{tabular}
\caption{Exact ordered $A$-orbit counts, second half of the triangular table.}
\label{tab:counts2}
\end{table}

Summing these entries gives $798644$. Exactly seven entries vanish.
Five are also excluded by the elementary finite or Euclidean triangle
quotients for $(2,3,3),(2,3,4),(2,3,5),(2,3,6),(2,4,4)$; the last two
are the attributed $(3,7)$ exclusion and the separately proved $(3,12)$
proper-cover exclusion. Orders one or two cannot supply an omitted
generating type: an order-one generator gives a cyclic group; if $R^2=1$
and $(RS)^2=1$, $R$ inverts $S$, so the group is dihedral. This proves
the complete rotation-type envelope in Theorem~\ref{thm:main}(1).

\section{Inverters and string triples}\label{sec:strings}
Suppose $(R,S)$ generates the centreless group $G$ and an automorphism
$b$ inverts both. Two such automorphisms have quotient fixing $R,S$, so
$b$ is unique. Its square fixes both, and $b\ne1$ because $p,q>2$;
therefore $b$ is an involution. Set
\[
 a=Rb,\qquad c=bS.
\]
Then $a,b,c$ are nonidentity involutions, $ab=R$, $bc=S$ and $ac=RS$.
Thus $a,c$ commute and are distinct. Conversely any such triple with
$\gen{ab,bc}=G$ supplies the simultaneous inverter $b$. The construction
is equivariant under $A$. Since $G$ has index two in $A$, all three
involutions have the same parity. Inner triples generate $G$; all-outer
triples generate $A$.

Here is the intersection-property argument, included to make the
polyhedron inference explicit. The cyclic groups $\gen R$ and $\gen S$
intersect trivially, since an element of their intersection commutes with
both generators and belongs to $Z(G)$. The adjacent dihedral groups are
$D_0=\gen{R,b}$ and $D_1=\gen{S,b}$. If their intersection were larger
than $\gen b$, comparison of forms $R^i b^\epsilon=S^j b^\delta$
would force a mixed-parity equality $R^i=S^j b$. For outer $b$ this
contradicts membership in $G$. For inner $b$, this common element is a
reflection in $D_1$: it centralizes $R$ and inverts $S$. Hence both
$(RS)^2=(RS^{-1})^2=1$. These relations make
$\gen{R^2,S^2}$ abelian and normal, with quotient elementary abelian
of order at most four: $R$ inverts $S^2$, and $S$ inverts $R^2$.
They would make $G$ solvable, a contradiction. Thus $D_0\cap D_1=\gen b$.
The other pair intersections with $\gen{a,c}\cong C_2^2$ are respectively
$\gen a$ and $\gen c$; otherwise the first equality would fail. These
are the full rank-three intersection property.

Conversely, an inner string triple generating the simple group has
rotation subgroup normal of index at most two; simplicity forces it to
be $G$. For an all-outer triple generating $A$, the rotation subgroup
lies in $G$ and has index at most two in $A$, so again equals $G$.
These observations recover the source's map/polyhedron assertions under
the specified classical interpretation.

Reversal sends $(a,b,c)$ to $(c,b,a)$ and rotations to
$(S^{-1},R^{-1})$. It preserves inner/outer parity. It is not an
additional identification in Tables~\ref{tab:counts1}--\ref{tab:counts2}.
If $d_p$ is the number of reversal-fixed inner orbits at $(p,p)$, the
number of inner strings after also identifying reversal is
\[
 \sum_{p<q}n_{\rm in}(p,q)+\sum_p
          \frac{n_{\rm in}(p,p)+d_p}{2}.
\]
Thus neither the source's triangular total nor a type count determines
the separate L026 rank-three total.

\section{Chirality and exact small splits}\label{sec:chirality}
There is one $A$-class of inner involutions and one of outer involutions.
For a representative $a_i$, $i=0,1$, let
$\kappa_m(i)$ count same-parity involutions $b$ with $|a_i b|=m$,
and put $h_i=|C_A(a_i)|$. Counting triples before imposing the second
adjacent order or generation bounds reflexible orbits by
\[
 n_i(p,q)\leq
 \left\lfloor\frac{\kappa_2(i)\min(\kappa_p(i),\kappa_q(i))}{h_i}\right\rfloor.
\]
The proof counts the choices of first involution, its partner and the
commuting endpoint, then divides by the free action on generating
triples. Reversal gives the minimum. Here
$h_0=322560$, $h_1=351120$; all $\kappa$ values are independently acquired
class-product counts in the named table $\texttt{ON.2}$.

Subtracting both bounds from an acquired generating lower bound proves
100 positive chiral types. The separately acquired literal certificates
prove another 23. Complete small-row arguments and the complete
$(6,6)$ split supply the remaining five. Their exact union is
$T\setminus\{(3,15)\}$, of size 128. A certificate here proves a
positive marking; an unfinished operation schedule is never used as an
exhaustion certificate. The complete exceptional split is
\[
 n(3,15)=6=5\text{ inner}+1\text{ outer}+0\text{ chiral}.
\]
The minimum-order row calculations give the following disjoint counts.
For order five the independently checked order-four cross-check is
excluded from this minimum-order sum.
\begin{center}
\begin{tabular}{rrrrr}\toprule
Minimum order&All&Inner&Outer&Chiral\\\midrule
3&284&39&15&230\\
4&5176&114&156&4906\\
5&7738&188&210&7340\\\bottomrule
\end{tabular}
\end{center}
The individually asserted small splits are
\begin{center}
\begin{tabular}{rrrrr}\toprule
Type&All&Inner&Outer&Chiral\\\midrule
$(3,10)$&7&0&1&6\\
$(4,5)$&18&0&2&16\\
$(4,6)$&43&0&1&42\\
$(5,5)$&26&2&2&22\\
$(6,6)$&165&8&1&156\\\bottomrule
\end{tabular}
\end{center}
These exact counts also settle the three isolated inner exclusions used
below. The accepted minimum-order-five row supplies an inner $(5,7)$
marking. The diagonal transfer in Section~\ref{sec:positives} supplies
the separately asserted inner $(7,7)$ marking.

\section{Complete negative type envelopes}\label{sec:envelopes}
The acquired inverse-class and fusion calculation gives two $G$-classes
of order 20 and two of order 31, with no inner inverter. The two
order-eight and four order-sixteen classes have no outer inverter.
These are necessary one-element conditions; no implication from separate
invertibility of two generators to a common inverter is used. Together
with the exact small zeros, every inner type lies in $I$.

There is one additional outer exclusion. If outer involutions $a,b$
satisfy $|ab|=6$, then $z=(ab)^3$ is an inner involution centralizing
$a,b$. Normalize $z$ and work in its acquired full centralizer of order
$322560$. The outer involutions there form one class of size $1344$.
For its representative, exactly 20 outer partners satisfy the normalized
order-six relation; they form one complete orbit with stabilizer order
12. Transporting the first member to the fixed outer involution $t$
puts the third member in the complete 1463-element domain from $J_1$.
The acquired full stabilizer partitions that domain into 141 orbits,
whose sizes sum to 1463. Their exact adjacent-product orders are
\[
 \{2,3,4,5,6,7,10,11,14,15,19,20,28,31\};
\]
order 12 never occurs. The order is computed on every orbit before any
generation filter. Thus no outer triple of type $(6,12)$ exists, and
reversal excludes the other orientation. Every outer type therefore
lies in $J$. This contradicts the particular outer existence assertion
on source PDF page 9 without excluding inner $(6,12)$ or changing the
total pair count for that type.

The upper envelopes have sizes $|I|=95$ and $|J|=97$. Their proof does
not require a complete census of all inversion orbits. Once a positive
certificate or a proved positive count exists at every member of each
envelope, equality follows. The exact positive union now equals these
two sets; its proof inputs are described next.

\section{Positive spectra and their final certificates}\label{sec:positives}
The positive union combines complete low-order rows, literal generating
witnesses, algebraic operations on accepted witnesses, and two counting
proofs. Their roles are different but all prove the same existential
predicate for one type and parity. No expected source count is an input
to a generation test.

For unit operations, a generating string $(a,b,c)$ with $[a,c]=1$
permits replacing the middle reflection by $a(ab)^e$ when
$\gcd(e,|ab|)=1$. The powers recover $ab$, hence recover $b$ and the
whole reflection subgroup. In the inner case that group is perfect, so
its even-word subgroup cannot be proper of index two; in the outer case
the even words generate the named index-two socle. Full word equalities,
orders and parity are checked for each submitted result. Classical
Petrie replacement $(a,b,c)\mapsto(ac,b,c)$ is used only on inner
strings: the generated group is unchanged, and perfectness again forces
the new rotations to generate it. Its attribution to the classical
Petrie operation is retained. The argument is not applied to mixed
parity outputs of an outer string.

The diagonal transfer uses an accepted string $(x,a,b)$ with
$|xa|=4$ and $[x,b]=1$. Put $c=a^x$. Then $a,c$ are distinct commuting
involutions, and $K=\gen{a,b,c}$ is normalized by $x$; it has index at
most two in the original reflection group. For an inner string,
perfectness gives $K=G$. For an all-outer string, $A'=G\leq K$ and
$a\notin G$ give $K=A$. Thus $(a,b,a^x)$ is a new generating string
with equal adjacent orders $|ab|$. The already proved intersection
property applies. This transfers the positive order-four rows to diagonal
types, including the required inner $(7,7)$ and $(10,10)$ and four
outer diagonal types. Subsequent fixed word calculations admit a positive
type only after independent literal reconstruction.

The counting proofs for inner $(6,10)$ and $(6,28)$ do not infer
generation merely from a free triple. For a selected triple with trivial
$A$-stabilizer, the centralizer $C$ of the rotation pair has
$C\cap C_A(b)=1$. Conjugation by the inner inverter $b$ is consequently
fixed-point-free on $C$, and the map
$u\mapsto u^{-1}u^b$ is bijective. Its image is inverted by $b$, so
$C$ is abelian of odd order. It lies in $G$ because $[A:G]=2$.
The acquired element centralizers of orders 72 and 20 (respectively 28)
then force $C=1$. Hence a free triple maps injectively to a pair before
any generation assumption. Only after this multiplicity issue is settled
are the raw maximal-overgroup bounds subtracted. The component proofs
give at least 12 inner $(6,10)$ orbits and at least two inner $(6,28)$
orbits; they do not identify a particular generating representative or
claim an exact count.

The last three complete pair-family acquisitions establish outer
$(10,12)$ with 12 ordered orbits, inner $(7,14)$ with four, and inner
$(7,10)$ with three. For the final remaining positive type, an acquired
inner $(4,11)$ string $(x,y,z)$ supplies a fixed first-order-eleven pair.
Write $s=yz$, $z=x_3^w$ and use the acquired transport $x_3^{h_3}=x$.
With right conjugation $u^v=v^{-1}uv$,
\[
 y^{s^6}=z,\qquad h=s^6w^{-1}h_3,\qquad y^h=x.
\]
The submitted third involution is the class-two representative conjugated
by the signed word $(1,2,-1,-1,-1,2,-1,-1,2)$ in the acquired centralizer
generators. Independent reconstruction checks the complete pair,
transport and rotation hashes. Direct full-array prime powers prove
the exact orders $11,7,2$ and the inversion equalities. A separately
recomputed complete point orbit has size $122760$. Every proper subgroup
is contained in a maximal subgroup, whereas none of the thirteen maximal
orders is divisible by $122760$. Orbit--stabilizer therefore proves
generation. This acquires one inner $(7,11)$ marking. It reuses the
named model and the generation predicate but checks the new orders by a
different direct-power path. It does not replay the 45 earlier negative
candidates or assert a complete count at this type.

The immutable positive union, including this final certificate, equals
$I$ and $J$. Combining it with Section~\ref{sec:envelopes} proves
Theorem~\ref{thm:main}(3)--(4). An interrupted word schedule was never
used to assert absence. The completed type spectra likewise assert no
uncomputed total number of reflexible markings.

\section{The source's constructive algorithm}\label{sec:algorithm}
Fix an order $p>2$. Choose one representative $s$ of every $G$-class
of order $p$, then one involution $\tau$ from every $\gen s$-orbit on
the involutions of $G$. Test
\[
 (r,s)=(\tau s^{-1},s),\qquad |r|\geq p,\qquad\gen{r,s}=G,
\]
and retain one representative under the full automorphism action.
Every required pair can first have its second entry conjugated to a
chosen $s$, and then its product involution moved to $\tau$ by a power
of $s$ without changing $s$. This proves coverage. The final full
automorphism comparison removes duplicates. Using the smaller cyclic
group rather than the whole centralizer affects redundancy, not
completeness; a cyclic orbit is not silently called a full centralizer
orbit.

The conjugation action on involutions is the action on right cosets of
$C_G(z)$ via $C_G(z)g\mapsto z^g$. It is faithful by simplicity and has
degree $g/161280=2857239$. Its points can retain their involution labels;
no unproved extraction from an abstract stabilizer is needed. The complete
roster also bounds every nonidentity centralizer by 161280: an element
centralizer lies inside the normalizer of its cyclic subgroup, and the
maximum normalizer among all roster rows of that order is an upper bound.
This uses all relevant rows rather than guessing which rows are cyclic.

Stopping when a known count is reached is valid only after checking that
every retained marking passes the tests and no two are equivalent. An
inclusion of finite sets with equal cardinality is equality. The stopping
count for $(20,31)$ must be $9118$, and the aggregate must be $798644$.
No runtime bound follows from this completeness proof. The historical
Magma runs and times reported in the source are not replayed here.

\section{Source corrections and limits}\label{sec:corrections}
\begin{center}\small
\begin{tabular}{p{.36\textwidth}p{.22\textwidth}p{.31\textwidth}}\toprule
Statement&Printed source&Acquired replacement\\\midrule
Table 1, $(20,31)$&7246&9118\\
Aggregate triangular count&796772&798644\\
$(3,15)$ reflexible split&4 inner, 2 outer&5 inner, 1 outer; no chiral\\
Outer type spectrum&98; includes $(6,12)$&97; excludes $(6,12)$\\
Adjacent $(3,12)$ prose&Both zero and positive&Zero; independent proper-cover proof\\\bottomrule
\end{tabular}
\end{center}
The accepted independent proof of the $(3,12)$ zero does not validate
the inadequate cited sample-pair argument. The source's complete table
and its mathematical claims are treated individually, including claims
that require correction. Historical computations remain historical;
the source's suggested future investigation of exponents and self-duality
is context, not an assertion of a completed census.

The rank-three reversal results derived along the way (114 inner and
156 all-outer reversal-fixed markings under the stated $A$ action) do
not settle L026's separate rank-three total or its higher-rank claims.
Likewise this treatment adds no Lean formalization, general conjecture
resolution or theorem-priority claim. Whole-source admission and public
registration are review actions separate from this mathematical assembly.

\section{Evidence, reproduction and source concordance}\label{sec:evidence}
The canonical text is this TeX file, its generated exact-table and
source-concordance inputs, and the companion structured binding list.
The assembly checker performs only finite integer/set and metadata joins:
it checks the 136 entries, unique source correction, aggregate, small
splits, 95/97 envelope equality and all 44 clause identifiers against
their exact acquired inputs. It does not relaunch a permutation census,
subgroup computation or witness search. The historical computational
proofs retain their own independently acquired receipts and input scopes.

The new binding list uses paths relative to the retained study root.
Publication must include the transitive canonical proof/code/input and
accepted-certificate closure, with an explicit resolver for absolute paths
inside unchanged historical records. A local synthesis check by itself
does not establish that those private files are already protected by a
public clone. Build scratch files are not mathematical prerequisites of
this manuscript. Historical records that themselves bind such files must
remain unaltered; a newly reviewed portable interface must distinguish
their audit history from the canonical reproduction inputs rather than
silently deleting a pinned premise.

\small
\begin{longtable}{p{.24\textwidth}rp{.45\textwidth}l}
\toprule Clause&PDF pp.&Disposition&Section\\\midrule\endhead
T1.1.1&3,7&recovered from accepted components&\ref{sec:scope}\\
T1.1.2&1,3,7&corrected complete&\ref{sec:scope}\\
T1.1.3&3,8&recovered complete chiral spectrum with corrected exception&\ref{sec:scope}\\
T1.1.4&3,8&recovered complete 95 inner types&\ref{sec:scope}\\
T1.1.5&3,8,9&corrected complete 97 outer types&\ref{sec:scope}\\
definitions&2,3,4,7&recovered or explicit classical premise&\ref{sec:strings}\\
T2.1&4&attributed borrowed theorem&\ref{sec:counting}\\
C2.2&5&recovered&\ref{sec:counting}\\
L2.3&5&recovered&\ref{sec:counting}\\
equation2.4&5&recovered&\ref{sec:counting}\\
algorithm3&5,6&recovered&\ref{sec:counting}\\
mobius&6&recovered general argument&\ref{sec:counting}\\
ONpremises&7,9&explicit borrowed inputs&\ref{sec:inputs}\\
p2excluded&7&elementary corollary&\ref{sec:counting}\\
table1&11&complete with one correction&\ref{sec:counting}\\
source\_prior\_results&7&attributed context&\ref{sec:corrections}\\
AlKadhi312&7&zero reproved and prose corrected&\ref{sec:corrections}\\
L5.1&8&corrected complete chiral spectrum and exception&\ref{sec:chirality}\\
L5.2&8&recovered&\ref{sec:envelopes}\\
L5.3&8&recovered&\ref{sec:envelopes}\\
split\_3\_10&8&recovered from complete row&\ref{sec:chirality}\\
split\_4\_5&8,9&recovered from complete row&\ref{sec:chirality}\\
split\_4\_6&8&recovered from complete row&\ref{sec:chirality}\\
split\_5\_5&8,9&recovered from complete row&\ref{sec:chirality}\\
inner\_witness\_5\_7&8&recovered from complete row&\ref{sec:positives}\\
inner\_witness\_7\_7&8&recovered diagonal transfer&\ref{sec:positives}\\
outer\_witness\_6\_12&9&refuted by complete accepted outer involution domain&\ref{sec:envelopes}\\
algorithm6&9&recovered with explicit completeness and corrected stopping total&\ref{sec:algorithm}\\
polyhedra&9&recovered with historical attribution&\ref{sec:strings}\\
row\_total\_3&9&recovered arithmetic&\ref{sec:chirality}\\
row\_chiral\_3&9,10&recovered from complete row&\ref{sec:chirality}\\
row\_inner\_3&9,10&recovered from complete row&\ref{sec:chirality}\\
row\_total\_4&9&recovered arithmetic&\ref{sec:chirality}\\
row\_chiral\_4&9,10&recovered from complete row&\ref{sec:chirality}\\
row\_inner\_4&9,10&recovered from complete row&\ref{sec:chirality}\\
row\_total\_5&9&recovered arithmetic&\ref{sec:chirality}\\
row\_chiral\_5&9,10&recovered from complete row&\ref{sec:chirality}\\
row\_inner\_5&9,10&recovered from complete row&\ref{sec:chirality}\\
historical\_timing&7,8,9,10&attributed historical observations&\ref{sec:corrections}\\
prevalence&10&qualitative context&\ref{sec:corrections}\\
outer\_row\_3&3,8,9&complete with correction&\ref{sec:chirality}\\
inner\_row\_3&3,8,9&complete with correction&\ref{sec:chirality}\\
complete\_row\_4&8,9&complete&\ref{sec:chirality}\\
complete\_row\_5&8,9,10&complete&\ref{sec:chirality}\\
\bottomrule\end{longtable}
\normalsize

\begin{thebibliography}{9}
\bibitem{CL} T. Connor and D. Leemans, Algorithmic enumeration of regular
maps, \emph{Ars Math. Contemp.} 10 (2016), 211--222.
\bibitem{Lattice} T. Connor and D. Leemans, An atlas of subgroup lattices
of finite almost simple groups, \emph{Ars Math. Contemp.} 8 (2015), 259--266.
\bibitem{Frob} G. Frobenius, Ueber Gruppencharaktere,
\emph{Berl. Ber.} (1896), 985--1021.
\bibitem{Hall} P. Hall, The Eulerian functions of a group,
\emph{Q. J. Math., Oxford Ser.} 7 (1936), 134--151.
\bibitem{Atlas} J. H. Conway, R. T. Curtis, S. P. Norton, R. A. Parker
and R. A. Wilson, \emph{Atlas of Finite Groups}, Oxford University Press, 1985.
\bibitem{CT} T. Breuer, \emph{The GAP Character Table Library}, retained
installed tables and exact acquisition interfaces.
\bibitem{Woldar} A. J. Woldar, On Hurwitz generation and genus actions of
sporadic groups, \emph{Illinois J. Math.} 33 (1989), 416--437.
\end{thebibliography}
\end{document}
